% !TeX root = ../navier-stokes-manuscript.tex
\subsection{More spatial regularity, not less}\label{sub:ThePressureCompleteMixedDerivativeTower}

Start with the base temporal rung \(V_0=u\).
Its critical-height energy is \(H=E_{1/2}\).
On any strict interval where the solution remains classical, write
\[
  E_s(t):=\frac12\norm{\Lambda^su(t)}_2^2,
  \qquad
  E_{s,\lambda}(t):=\frac12\norm{\Lambda^su_\lambda(t)}_2^2,
\]
for \(s\geq0\).
The same \NavierStokes{} scaling calculation in \Cref{prop:CriticalSobolevScaling} gives
\begin{equation}\label{eq:BodySpatialEnergyScaling}
  E_{s,\lambda}(t)
  =\lambda^{2s-1}E_s(\lambda^2t).
\end{equation}
At \(s=\tfrac12\), the exponent vanishes:
\[
  H_\lambda(t)=H(\lambda^2t).
\]
Differentiate this identity \(n\) times:
\begin{equation}\label{eq:BodyCriticalHeightTemporalScaling}
  H_\lambda^{(n)}(t)
  =\lambda^{2n}H^{(n)}(\lambda^2t).
\end{equation}
The spatial energy at height \(n+\tfrac12\) carries exactly the same factor:
\begin{equation}\label{eq:BodySpatialHeightScalingDiagonal}
  E_{n+1/2,\lambda}(t)
  =\lambda^{2n}E_{n+1/2}(\lambda^2t).
\end{equation}
Equal scale weight does not identify or compare the two quantities.
It only places the spatial energy \(E_{n+1/2}\) beside the critical response \(H^{(n)}\).
To see why that spatial height enters, return to the equation.

Applying \(\Lambda^s\) to Navier--Stokes and pairing the result with \(\Lambda^su\) gives the viscous contribution
\[
  \nu\pair{\Lambda^su}{\Lambda^s\Delta u}_{L^2}
  =-\nu\norm{\Lambda^{s+1}u}_2^2
  =-2\nu E_{s+1}.
\]
Since the global pressure pairing vanishes against the divergence-free velocity, define the nonlinear production at height \(s\) by
\[
  \mathcal P_s
  :=-
  \operatorname{Re}
  \pair{\Lambda^su}
  {\Lambda^s\bigl((u\cdot\nabla)u\bigr)}_{L^2}.
\]
The complete spatial-energy row is
\begin{equation}\label{eq:BodyCriticalEnergyBalance}
  E_s'=\mathcal P_s-2\nu E_{s+1}.
\end{equation}

At critical height, one time derivative reaches the next spatial energy:
\[
  H'=\mathcal P_{1/2}-2\nu E_{3/2}.
\]
Differentiate again and use the same row at \(s=\tfrac32\):
\[
\begin{aligned}
  H''
  &=\partial_t\mathcal P_{1/2}-2\nu E_{3/2}'
  \\
  &=\partial_t\mathcal P_{1/2}
  -2\nu\mathcal P_{3/2}
  +(2\nu)^2E_{5/2}.
\end{aligned}
\]
The third derivative reaches one height farther:
\[
  H'''
  =\partial_t^2\mathcal P_{1/2}
  -2\nu\partial_t\mathcal P_{3/2}
  +(2\nu)^2\mathcal P_{5/2}
  -(2\nu)^3E_{7/2}.
\]
Each differentiation keeps every transfer term already produced and exposes one higher viscous energy.
For \(n\geq1\), induction gives
\begin{equation}\label{eq:BodyCriticalHeightSpatialAscent}
  H^{(n)}
  =(-2\nu)^nE_{n+1/2}
  +\sum_{j=0}^{n-1}
  (-2\nu)^j
  \partial_t^{\,n-1-j}\mathcal P_{j+1/2}.
\end{equation}
Because \(\nu>0\), the top spatial energy can be isolated:
\begin{equation}\label{eq:BodyCompletedCriticalRow}
  E_{n+1/2}
  =(-2\nu)^{-n}
  \left(
    H^{(n)}
    -\sum_{j=0}^{n-1}
      (-2\nu)^j
      \partial_t^{\,n-1-j}\mathcal P_{j+1/2}
  \right).
\end{equation}
The induction in \Cref{prop:CriticalHeightSpatialAscent} isolates the higher Sobolev energy, but does not bound it.
A bound requires control of the complete right-hand side of \eqref{eq:BodyCompletedCriticalRow}.

The scalar response \(H^{(n)}\) is itself built from the temporal rows.
Since
\[
  H(t)
  =\frac12
  \pair{\Lambda^{1/2}V_0(t)}
       {\Lambda^{1/2}V_0(t)}_{L^2},
\]
Leibniz's rule gives
\begin{equation}\label{eq:BodyHeightDerivativeTemporalContraction}
  H^{(n)}(t)
  =\frac12\sum_{a=0}^{n}\binom{n}{a}
  \operatorname{Re}
  \pair{\Lambda^{1/2}V_a(t)}
       {\Lambda^{1/2}V_{n-a}(t)}_{L^2}.
\end{equation}
The three entries on one rung remain different objects.
The field response is \(V_n\), the scalar response \(H^{(n)}\) contracts pairs of temporal rows, and the viscous calculation exposes the spatial energy \(E_{n+1/2}\).

\Needspace{10\baselineskip}
\begin{center}
\captionsetup{hypcap=false}
\captionof{table}{The first time--space rungs of the critical-height calculation.}
\label{tab:BodyTemporalSpatialDiscovery}
\small
\renewcommand{\arraystretch}{1.2}
\begin{tabular}{@{}c c c c@{}}
\toprule
order & Eulerian response & critical response & spatial height \\
\midrule
\(0\) & \(V_0=u\) & \(H=E_{1/2}\) & \(E_{1/2}\) \\
\(1\) & \(V_1=u_t\) & \(H'\) & \(E_{3/2}\) \\
\(2\) & \(V_2=u_{tt}\) & \(H''\) & \(E_{5/2}\) \\
\(3\) & \(V_3=u_{ttt}\) & \(H'''\) & \(E_{7/2}\) \\
\(n\) & \(V_n=\partial_t^nu\) & \(H^{(n)}\) & \(E_{n+1/2}\) \\
\bottomrule
\end{tabular}
\end{center}

The temporal column records successive Eulerian responses of the same preterminal velocity field.
The spatial column records the higher Sobolev energies exposed beside them:
\(E_{1/2},E_{3/2},E_{5/2},E_{7/2},\ldots\).

This is one exposed ladder, not yet a regularity estimate.
The transfer terms can still carry the size of its highest rung, and the calculation has so far opened the ladder only inside the original velocity response \(V_0=u\).

The same calculation can be made inside every temporal response.
For \(r\in\Nzero\), let \(E_{r,s}:=\tfrac12\norm{\Lambda^sV_r}_2^2\), and define the complete transport--pressure transfer in the \(r\)-th recurrence by
\[
  \mathcal P_{r,s}
  :=-
  \operatorname{Re}
  \pair{\Lambda^sV_r}
  {\Lambda^s\!\left[
    \sum_{a=0}^{r}\binom{r}{a}(V_a\cdot\nabla)V_{r-a}
    +\nabla Q_r
  \right]}_{L^2}.
\]
Pairing the \(r\)-th temporal equation with \(\Lambda^{2s}V_r\) gives
\begin{equation}\label{eq:BodyEveryTemporalRowSpatialEnergy}
  E_{r,s}'=\mathcal P_{r,s}-2\nu E_{r,s+1}.
\end{equation}
After any finite ascent depth \(k\geq1\),
\begin{equation}\label{eq:BodyEveryTemporalRowCompletedHeight}
  \partial_t^kE_{r,1/2}
  =(-2\nu)^kE_{r,k+1/2}
  +\sum_{j=0}^{k-1}
  (-2\nu)^j
  \partial_t^{\,k-1-j}\mathcal P_{r,j+1/2}.
\end{equation}
The index \(r\) chooses the temporal derivative order, while \(k\) chooses how high the spatial ladder has been followed inside that response.
These are independent coordinates of the same mixed jet.

One spatial ladder inside every temporal response forms the grid
{\small
\[
\begin{array}{c|ccccc}
 & H^0_x & H^1_x & H^2_x & H^3_x & \cdots\\
\hline
V_0=u      & \bullet & \bullet & \bullet & \bullet & \cdots\\
V_1=u_t    & \bullet & \bullet & \bullet & \bullet & \cdots\\
V_2=u_{tt} & \bullet & \bullet & \bullet & \bullet & \cdots\\
V_3=u_{ttt}& \bullet & \bullet & \bullet & \bullet & \cdots\\
\vdots     & \vdots  & \vdots  & \vdots  & \vdots  & \ddots
\end{array}
\]
}
The bullets mark derivatives of the same velocity field, not bounds already proved for them.
Moving down chooses a later Eulerian response, and moving across chooses a higher Sobolev height inside that response.
Fix \(T<T_*\) and set \(I=(0,T)\).
Control of every coordinate on this strict classical interval would place the velocity in the mixed intersection
\begin{equation}\label{eq:BodyDesiredMixedIntersection}
  u\in
  \bigcap_{r=0}^{\infty}
  \bigcap_{m=0}^{\infty}
  H^r(I;H^m(\R^3)).
\end{equation}

The grid is infinite, but any chosen temporal order and spatial height lie inside a finite block.
Fix \(N\) and \(M\); the rows through \(N\) and the columns through \(M\) form one rectangle.
For each fixed rectangle, the datum-generated whole-terminal estimate proved in \Cref{thm:WholeTerminalCompatibleRectangleEstimate} controls the entire rectangle uniformly over \(T<T_*\); its bound may depend on \(N\) and \(M\).
Every coordinate appears in some rectangle, and the rectangles agree wherever they overlap.
That agreement allows Section~\ref{sec:TheIntersectionInsideTheTower} to pass from the finite rectangles to their compatible intersection.
